\documentclass[12pt]{article}
\usepackage[OT4]{fontenc}
\usepackage[polish]{babel}
\usepackage[legalpaper, margin=2cm]{geometry}
\usepackage{mathtools}
\usepackage{float}
\usepackage{hyperref}
\usepackage{svg}
\usepackage{graphicx}
\usepackage{placeins}
\usepackage{chngcntr}

\newif\ifpokazsekcjaszesc
\pokazsekcjaszesctrue

\newcommand{\wykrespngwniosek}[4]{%
\begin{figure}[H]
    \centering
    \ifpokazsekcjaszesc
        \includegraphics[width=0.72\textwidth,height=0.24\textheight,keepaspectratio]{#1}%
    \else
        \fbox{%
            \parbox[c][0.20\textheight][c]{0.72\textwidth}{%
                \centering
                \texttt{#1}\\[0.4em]
                \small miejsce na wykres
            }%
        }%
    \fi
    \caption{#2}
    \label{fig:#3}
\end{figure}
\noindent\textbf{Wniosek:} #4
\par\vspace{0.2cm}
}


\title{Sprawozdanie z drugiej połowy projektu\\
Przedmiotu "Wizualizacja Danych"}
\author{Grupa Pon 14:15 - 17:00, \\zespół 4 - Jan Domański, Maciej Kurzydłowski, Jan Machowski}
\date{sem. 26L}

\begin{document}
\maketitle
\tableofcontents
\newpage

% Proponowana struktura rozdziału
% 1. Źródło danych
% 2. Opis zbioru i istotnych pojęć - o czym jest, pojęcia z nim związane, kolumny, czy dane są kateg / liczb/
% 3. Wizualizacja + krótki opis co pokazuje
% 4. Wnioski


% \subsubsection{}
% \begin{figure}[H]
%     \includegraphics[width=\textwidth]{phts/}
%     \caption{}
% \end{figure}


\section{Trzęsienia ziemi na świecie}

Przygotowane przez: Maciej Kurzydłowski

\subsection{Źródło danych}
\href{}{}\\
\url{}

\subsection{Opis zbioru i istotnych pojęć}
\begin{itemize}
    \item 
\end{itemize}


\subsection{Wizualizacje zbioru}
\subsubsection{}
\begin{figure}[H]
    \includesvg[width=\textwidth]{phts/1/svg/}
    \caption{}
\end{figure}

\newpage
\subsection{Wnioski}
\begin{enumerate}
    \item 
\end{enumerate}

\end{document}